%% Macros de dados do documento
\autor{Artur Dias Rodrigues Vicente}
\titulo{Supervisão e controle de nível com lógica Fuzzy: Uma integração baseada em OPC}
\instituicao{Instituto Federal do Espírito Santo}
\campus{Campus Serra}
\curso{ENGENHARIA DE CONTROLE E AUTOMAÇÃO}
\local{Serra -- ES}
\data{2025}


\preambulo {Trabalho de conclusão de curso, apresentado como parte das atividades para obtenção do título de Bacharel em Engenharia de Controle e Automação, do curso de Engenharia de Controle e Automação do Instituto Federal do Espírito Santo.
\vspace{0.5cm}\\ Orientador: Prof. Dr. Daniel Cruz Cavalieri
}

\examinadori{Prof. Dr. Daniel Cruz Cavalieri}{Instituto Federal do Espírito Santo}
{Campus Serra}
\examinadorii{Prof. Me. Rafael Emerick Zape de Oliveira}{Instituto Federal do Espírito Santo}
{Campus Serra}
\examinadoriii{Prof. Dr. Flávio Garcia Pereira}{Instituto Federal do Espírito Santo}
{Campus Serra}

\approvaldate{28}{Agosto}{2025}


% Ficha catalográfica
\fichabiblioteca{
    Dados Internacionais de Catalogação-na-Publicação (CIP) \\
    (Biblioteca Nilo Peçanha do Instituto Federal do Espírito Santo)
}
\fichaautor{Elaborada por XXXXXXXXXXXXXXXXXXXXX – CRB-X/ES - XXX}
\fichatags{
    1. Redes neurais (Computação).  2. Expressão facial – Processamento de dados. 3. Sistemas de reconhecimento de padrões. 4. Processamento de imagens – Técnicas digitais. 5. Percepção de padrões. 6. Engenharia Elétrica. I. XXXXXXXX, XXXXXXXXXX XXXXXXXXXX.  II. Instituto Federal do Espírito Santo. III. Título.
}

\cutter{X999y}
\cdd{000.00}
\selectlanguage{brazil}

% Arial (para usar times-new-roman, comente as três linhas abaixo)
\usepackage{helvet}
\renewcommand{\familydefault}{\sfdefault}
\renewcommand{\ABNTEXchapterfont}{\bfseries}
% Comandos para revisar o texto
\usepackage{easyReview} 


%% Elementos opcionais
\editardedicatoria{\input{pre_textuais/dedicatoria}}
\editaragradecimentos{Agradeço a Deus por me conceder força, sabedoria e perseverança ao longo desta jornada. Sou profundamente grato à minha família pelo apoio incondicional, incentivo e amor, fundamentais para a realização deste trabalho. 

Estendo meus sinceros agradecimentos a todos os professores que me orientaram com conhecimento, paciência e dedicação, contribuindo significativamente para meu crescimento acadêmico e pessoal.

Por fim, reconheço o Instituto Federal do Espírito Santo por proporcionar recursos, oportunidades e um ambiente propício ao aprendizado e desenvolvimento.}
\editarepigrafe{\thispagestyle{empty}

\begin{minipage}{1cm}

\end{minipage}

\vfill

\begin{center}
\begin{minipage}{11cm}
"A ciência pode divertir e fascinar, mas é a engenharia que muda o mundo."
\end{minipage}
\end{center}

\vspace{1.5cm}

\begin{flushright}
(Asimov, 1988)
\end{flushright}}
